% One-time imported draft. Editable native source; no Markdown regeneration.
\section{Step 2 Declare the
  messages}\label{001-console-calculation-s02-step-2-declare-the-messages}

The .data directive selects the data section. A label gives an address a
name. The .asciiz directive emits the text bytes followed by a zero
byte; the string-printing service uses that zero to find the end. The
colon belongs after the label name.

\subsection{Write or
  inspect}\label{001-console-calculation-s02-write-or-inspect}

Add these declarations inside the student data region. Preserve spaces
at the ends of the two prompts.

\MARSCodeLabel{MIPS assembly}
\begin{minted}{mars-idea}
prompt:  .asciiz "Enter a small integer: "
answer:  .asciiz "Twice your number plus one is: "
newline: .asciiz "\n"
\end{minted}

\subsection{Expected
  observation}\label{001-console-calculation-s02-expected-observation}

The file should still assemble and exit silently. Declaring data does
not print it.

\subsection{Check your
  understanding}\label{001-console-calculation-s02-check-your-understanding}

What does the label prompt name?

\begin{itemize}
  \tightlist
  \item
        \textbf{A} The number of times the prompt will print.
  \item
        \textbf{B} A register dedicated to console input.
  \item
        \textbf{C} The address of the first byte of the prompt string.
\end{itemize}
